\section{访谈折射出的创业方法论}

\subsection{为什么模型公司不能只做模型}
整场访谈虽然围绕 K2、Muon、RL 和 Agent 展开，但杨植麟反复强调的其实是一个更大的判断：\textbf{模型能力、评估体系和产品环境必须同步演进}。如果只把公司理解为 ``训练一个更强基础模型''，就会低估后续最困难的环节，因为真正把能力变成用户价值的，是模型外面的系统。

\begin{importantbox}{从能力供给到任务闭环，价值链正在后移}
\begin{itemize}
    \item 预训练和后训练决定模型的能力上限；
    \item 评估体系决定这些能力是否能被稳定识别和持续迭代；
    \item 产品环境决定模型是否能在真实工作流中把能力兑现出来。
\end{itemize}
这也是为什么访谈里会同时出现 K2、Agent、Benchmark 和一方产品等话题。它们不是分散的热点，而是同一条价值链上的不同环节。
\end{importantbox}

杨植麟没有把 ``产品'' 理解为最后包一层聊天界面，而是更接近 ``给模型搭一个可以长期学习和执行的环境''。在这种视角下，产品既是分发渠道，也是训练环境的一部分。模型公司如果忽视这一点，很容易在 demo 层面看起来很强，但在真实任务上无法建立复利。

\subsection{创业公司的节奏与窗口}
这场访谈也隐含了一种创业节奏判断：在范式剧烈变化期，团队需要优先抓那些能够形成复利的基础能力，而不是追逐短期热点。无论是 Muon 优化器、Token Efficiency、Agent 泛化，还是评估体系，本质上都在争夺未来几年的技术坡度。

\begin{knowledgebox}{访谈中可以抽象出的三层窗口}
\begin{enumerate}
    \item \textbf{算法窗口}：谁能率先找到更高 Token Efficiency、更稳定 RL 和更好的 verifier，谁就能以更低成本扩展能力；
    \item \textbf{系统窗口}：谁能把工具、环境、记忆和评估接到一起，谁就更容易把模型能力沉淀为产品体验；
    \item \textbf{组织窗口}：谁能让研究、工程、产品数据形成快反馈回路，谁就更容易持续迭代。
\end{enumerate}
\end{knowledgebox}

\begin{warningbox}{最危险的误判是把短期领先误认为长期壁垒}
访谈对 Benchmark 的警惕同样适用于公司判断。短期的演示效果、单点 benchmark 提升，甚至某次产品爆发，都不必然构成长期壁垒。真正的壁垒来自：能否持续吸收高质量数据、能否建立可靠评估、能否把用户任务反哺到训练与产品闭环之中。
\end{warningbox}

\textit{``问题是不可避免的，问题也是可以解决的。''} 如果把这句话放到创业语境里，它意味着团队不应该期待某个 ``终局模型'' 一次性解决所有问题，而应该围绕关键瓶颈建立一套持续爬坡的机制。K2 只是其中一个阶段性的高点，而不是路线结束。

\subsection{本章小结}
从这场访谈能够抽象出的创业方法论是：不要把模型、评估和产品分开看，而要把它们视作同一个系统的不同层。未来能赢下长期竞争的，不只是训练出强模型的团队，而是能把强模型持续变成真实任务闭环的团队。


